\begin{figure}[t]
\centering
\resizebox{\textwidth}{!}{%
\begin{tikzpicture}[
  node distance=6mm and 7mm,
  box/.style={draw,rounded corners=2pt,minimum height=9mm,align=center,font=\small,inner xsep=5pt},
  flow/.style={-{Stealth[length=2mm]},thick,prgray},
  note/.style={font=\scriptsize,align=center,text=prgray}
]
\node[box,fill=prlightgold] (choice) {ordinary tree\\or controlled motif};
\node[box,fill=prlightblue,right=of choice] (tree) {canonical\\process tree};
\node[box,fill=prlightteal,right=of tree] (nets) {two Petri\\realizations};
\node[box,fill=prlightred,right=of nets] (cert) {exact / bounded /\\isomorphic evidence};
\node[box,fill=prlightblue,below=10mm of nets] (pool) {certified or sampled\\trace pool};
\node[box,fill=prlightteal,left=of pool] (views) {two clean views\\128 traces each};
\node[box,fill=prlightgold,left=of views] (rows) {four aligned rows\\per family};

\draw[flow] (choice) -- (tree);
\draw[flow] (tree) -- (nets);
\draw[flow] (nets) -- (cert);
\draw[flow] (cert) |- (pool);
\draw[flow] (nets) -- (pool);
\draw[flow] (pool) -- (views);
\draw[flow] (views) -- (rows);
\draw[flow] (tree) |- (rows);
\node[note,below=1mm of rows] {tree + log + marked graph + IDs + 128-D signature};
\end{tikzpicture}
}
\caption{Family-level generation. Family grouping is established before flattening, so all alternative realizations and observations remain in one data split.}
\label{fig:corpus}
\end{figure}
